\bookchapter{Глава 3: Подключаем телефон к «линии жизни» — Сеть}

\chapterintro{Без интернета смартфон «глупеет» наполовину. Сеть — это линия жизни современного телефона. Понять её, правильно использовать и защищать — обязательный навык для каждого пользователя.}

\booksection{Не теряйся в сети — различаем «дороги» и «машины» Wi-Fi и мобильного трафика}

У многих людей в верхней части экрана одновременно горят значки Wi-Fi и сигнала, но суть различий между ними остаётся загадкой. Классическое заблуждение:\textbf{подключён к Wi-Fi, но страницы не грузятся}— потому что роутер отключён от интернета, а телефон автоматически не переключился на мобильный трафик. Чтобы разобраться в этом до конца, построим модель «дорога и машина».

\begin{tcbinfo}{Ключевая метафора: дорога и машина}
Представьте интернет как город.\textbf{«Дорога»}— это канал между вашим телефоном и интернетом,\textbf{«машина»}— данные, которые по ней едут. Wi-Fi и мобильный трафик — две разные дороги, но данные (машины) сами по себе одинаковы.

\end{tcbinfo}

\booksubsection{Суть различий: локальный канал против удалённой магистрали}

Wi-Fi по своей природе — это\textbf{технология локальной сети}, как беспроводная дорога, проложенная у вас под домом, которая передаёт данные от роутера к телефону. Но если не оплачен широкополосный доступ или оборван оптоволоконный кабель — дорога есть, а «автомагистраль» наружу отрезана, и вы никуда не попадёте. А мобильный трафик (мобильные данные) передаётся через вышки 4G/5G оператора — это как выезд прямо на скоростную линию оператора: есть сигнал — есть прямой выход в интернет.

\booksubsection{Стоимость и сценарии: фиксированная абонплата против оплаты по объёму}

Wi-Fi — обычно разовая инвестиция (ежемесячная плата за широкополосный доступ), пока оплачено электричество — вся семья пользуется безлимитно. Идеально для скачивания фильмов, игр и других трафикоёмких задач. Мобильный трафик тарифицируется по объёму, при превышении пакета скорость может упасть или спишется дополнительная плата. Лучше подходит для лёгких задач в метро, командировках и других мобильных сценариях.

\begin{tcbpractice}{Попробуйте: проверьте свой тариф}
Откройте [Настройки] > [Мобильная сеть] > [Управление трафиком] и посмотрите, сколько трафика осталось на этот месяц и какие приложения расходуют больше всего. Когда знаешь цифры — выходишь из дома спокойно.

\end{tcbpractice}

\booksubsection{Безопасность: публичный Wi-Fi — скрытое «место для фишинга»}

С точки зрения безопасности мобильные данные обычно надёжнее публичного Wi-Fi. Мобильные данные шифруются оператором, а публичный Wi-Fi действительно легко превратить в «поддельную точку доступа» для перехвата незашифрованного трафика. Но здесь важно различать: при повсеместном распространении HTTPS злоумышленник может перехватить\textbf{какие сайты вы посещаете}(доменные имена) и\textbf{незашифрованный HTTP-контент}; а сообщения в WeChat, пароли от банковских приложений и другие данные, защищённые HTTPS, злоумышленник\textbf{не может прочитать напрямую}— если только не проведёт более сложную атаку SSL-stripping, а пользователь проигнорирует предупреждение сертификата. Но даже в этом случае, совершая платёжные операции вне дома, безопаснее всего использовать мобильный трафик.

\textbf{Напоминание о безопасности}\par

В аэропортах, торговых центрах, кафе и других общественных местах не подключайтесь к подозрительным бесплатным сетям Wi-Fi (например, «Free\_WiFi\_2024»). Даже если подключились — ни в коем случае не входите в банковские приложения и не вводите пароли.

\booksubsection{Как определить на телефоне: читаем «шифр» строки состояния}

Как узнать, через что телефон сейчас выходит в интернет? На примере Huawei: если рядом с сигналом оператора или сигналом WLAN в строке состояния появляется\textbf{значок двусторонней стрелки}— значит, интернет идёт через этот канал; если горят оба одновременно — работает ускорение сети, телефон использует оба канала сразу.

У iPhone всё проще: горит значок Wi-Fi — значит, используется Wi-Fi; не горит — смотрите на деления сигнала в правом верхнем углу, это мобильный трафик. Некоторые Android-телефоны также показывают в строке состояния «5G» или «4G», прямо сообщая текущий тип мобильной сети.

\booksection{Мобильная сеть (4G/5G) и трафик}

Мобильная сеть — основа подключения к интернету в любое время и в любом месте. От 2G до 5G скорость выросла в тысячи раз, но принцип примерно тот же.

\booksubsection{Стандарты сети:}

\begin{itemize}
  \item \textbf{4G LTE}: самое широкое покрытие на сегодня, скорость достаточная (теоретически до 100 Мбит/с+), относительно низкое энергопотребление.
  \item \textbf{5G}: высокая скорость (теоретически до 10 Гбит/с), но покрытие ещё расширяется, высокое энергопотребление. Подходит для скачивания больших файлов, просмотра HD-видео.
  \item \textbf{5G SA против NSA}: SA — автономная сеть, NSA — с опорой на базовые станции 4G. SA даёт лучший опыт, но требует поддержки и оператором, и телефоном.
\end{itemize}

\booksubsection{Управление трафиком:}

\begin{itemize}
  \item \textbf{Что такое трафик}: это объём данных, переданных телефоном при выходе в интернет. Видео и картинки расходуют больше всего трафика, текст и голос — относительно экономно.
  \item \textbf{Советы по экономии}: включите в настройках «Режим экономии трафика», обновляйте приложения и смотрите HD-видео только по Wi-Fi.
  \item \textbf{Деления сигнала $\neq$ скорость}: сильный сигнал означает лишь хорошее соединение, скорость зависит от загруженности базовой станции и расстояния до неё.
\end{itemize}

\textbf{Важное предупреждение: отключите [Роуминг данных] перед выездом за границу!}\par

Использование мобильных данных за границей может привести к астрономическим счетам! Перед выездом обязательно отключите роуминг данных в настройках или заранее приобретите местную SIM-карту / международный роуминговый пакет.

\booksection{Wi-Fi: подключение, общий доступ и безопасность}

Wi-Fi — основной способ выхода в интернет дома и в офисе: высокая скорость, не расходует трафик. Но Wi-Fi постоянно развивается, и различия между стандартами весьма значительны.

\booksubsection{Подключение к Wi-Fi:}

Зайдите в [Настройки] > [WLAN], выберите сеть и введите пароль. Сейчас многие телефоны поддерживают вход по QR-коду, что ещё удобнее.

\booksubsection{Поделиться паролем:}

И Android, и iOS позволяют быстро поделиться паролем с другом с уже подключённого телефона. Android генерирует QR-код, iOS может передать пароль автоматически при приближении.

\booksubsection{Эволюция стандартов Wi-Fi:}

Стандарты Wi-Fi разрабатываются IEEE, каждое поколение — значительный шаг вперёд:

\begin{center}
\begin{tabularx}{\linewidth}{|X|X|X|X|X|}
\hline
\textbf{Стандарт} & \textbf{Протокол} & \textbf{Диапазон} & \textbf{Макс. скорость} & \textbf{Особенности} \tabularnewline
\hline
Wi-Fi 4 & 802.11n & 2.4/5 ГГц & 600 Мбит/с & Устаревший стандарт, покрытие приемлемое \tabularnewline
Wi-Fi 5 & 802.11ac & 5 ГГц & 3,5 Гбит/с & Самый распространённый, скорость достаточная \tabularnewline
Wi-Fi 6 & 802.11ax & 2.4/5 ГГц & 9,6 Гбит/с & Отличная работа с множеством устройств, низкая задержка \tabularnewline
Wi-Fi 6E & 802.11ax & 2.4/5/6 ГГц & 9,6 Гбит/с & Добавлен диапазон 6 ГГц, без перегрузок \tabularnewline
Wi-Fi 7 & 802.11be & 2.4/5/6 ГГц & 46 Гбит/с & Сверхскоростная передача, агрегация нескольких каналов \tabularnewline
\hline
\end{tabularx}
\end{center}

\begin{itemize}
  \item \textbf{Ключевые технологии Wi-Fi 6}: OFDMA (множественный доступ с ортогональным частотным разделением) позволяет нескольким устройствам передавать данные одновременно, MU-MIMO (многопользовательский MIMO) повышает эффективность параллельной работы устройств, TWT (целевое время пробуждения) экономит заряд батареи.
  \item \textbf{Wi-Fi 6E}: добавляет к Wi-Fi 6 диапазон 6 ГГц, предоставляя больше каналов и снижая помехи. Требуется поддержка и телефоном, и роутером.
  \item \textbf{Wi-Fi 7}: внедряет MLO (многоканальную работу), позволяя одновременно использовать несколько диапазонов для передачи данных, что значительно снижает задержку. Сверхширокий канал 320 МГц обеспечивает скачок скорости. Флагманские чипсеты, такие как Snapdragon 8 Elite Gen5 и Dimensity 9400, уже поддерживают Wi-Fi 7; модели Xiaomi 17, Huawei Mate 80, Samsung S26 могут использовать сеть Wi-Fi 7 (требуется роутер с Wi-Fi 7).
  \item \textbf{2.4 ГГц против 5 ГГц}: 2.4 ГГц лучше проходит сквозь стены, но медленнее и с большими помехами; 5 ГГц быстрее, но хуже проходит сквозь стены. Диапазон 6 ГГц в Wi-Fi 6E/7 самый быстрый, но с наихудшей проникающей способностью — подходит для использования в одной комнате.
\end{itemize}

\booksubsection{Риски публичного Wi-Fi:}

В публичных сетях Wi-Fi хакеры могут перехватывать данные или создавать фишинговые точки доступа. Правила:

\begin{itemize}
  \item Не входите в банковские аккаунты и не вводите пароли.
  \item Лучше использовать VPN для шифрования соединения.
  \item Не подключайтесь автоматически к незнакомым сетям Wi-Fi.
  \item Пользуйтесь сайтами HTTPS (значок замка в адресной строке).
\end{itemize}

\booksection{Bluetooth: универсальный боец ближней беспроводной связи}

Bluetooth — возможно, одна из самых часто используемых беспроводных технологий на вашем телефоне: беспроводные наушники, умные часы, автомобильные системы, клавиатуры и мыши — всё работает через него.

\booksubsection{Технический принцип:}

Bluetooth работает в диапазоне ISM 2.4 ГГц, используя технологию расширения спектра со скачкообразной перестройкой частоты (FHSS), распределяя данные по 79 каналам с быстрым переключением для эффективного избегания помех. Дальность связи Bluetooth делится на три класса: Class 1 — около 100 м, Class 2 — около 10 м (обычно используется в телефонах), Class 3 — около 1 м.

\booksubsection{Эволюция версий:}

\begin{itemize}
  \item \textbf{Bluetooth 4.0}: внедрение Bluetooth Low Energy (BLE), позволившее устройствам с маленькой батареей, таким как браслеты и пульсометры, работать без проводов — автономность выросла с дней до месяцев.
  \item \textbf{Bluetooth 4.2}: повышение скорости передачи и защиты конфиденциальности, поддержка прямого подключения IPv6.
  \item \textbf{Bluetooth 5.0}: удвоение скорости передачи (2 Мбит/с), увеличение дальности в 4 раза (теоретически 300 м), увеличение вещательной способности в 8 раз. Поддержка Dual Audio (одновременное подключение двух пар наушников).
  \item \textbf{Bluetooth 5.1}: внедрение функции определения направления (AoA/AoD), точность позиционирования до сантиметров.
  \item \textbf{Bluetooth 5.2}: новый стандарт LE Audio, поддержка кодека LC3 (лучшее качество звука, ниже энергопотребление), многопоточное аудио, широковещательное аудио.
  \item \textbf{Bluetooth 5.3/5.4}: дальнейшая оптимизация энергопотребления и стабильности соединения, улучшение периодического вещания.
\end{itemize}

\textbf{Распространение LE Audio}: в 2025–2026 годах LE Audio вступил в фазу быстрого распространения. Новые TWS-наушники основных брендов (такие как Huawei FreeBuds Pro 5, Xiaomi Buds 6 Pro, Samsung Galaxy Buds3 Pro и др.) уже поддерживают LE Audio и кодек LC3, обеспечивая меньшую задержку и лучшее качество звука. При покупке новых наушников рекомендуется отдавать предпочтение моделям с поддержкой LE Audio.

\booksubsection{Сценарии применения:}

\begin{itemize}
  \item \textbf{Передача аудио}: беспроводные наушники, автомобильная акустика. LE Audio в Bluetooth 5.2 даёт качественный скачок в звуке и задержке.
  \item \textbf{Передача данных}: передача файлов и фото на близком расстоянии. Медленнее, чем Wi-Fi Direct, но выигрывает в простоте.
  \item \textbf{Сопряжение устройств}: умные часы, браслеты, клавиатуры и мыши, игровые контроллеры.
  \item \textbf{Сервисы геолокации}: Bluetooth-маячки (Beacon) для навигации внутри помещений, отслеживания предметов (например, Apple AirTag).
  \item \textbf{Умный дом}: ячеистая сеть Bluetooth Mesh для управления лампами, замками и т.д.
\end{itemize}

\booksection{ИК-пульт: телефон как универсальный пульт ДУ}

Инфракрасное дистанционное управление — одна из старейших технологий беспроводного управления, но на телефоне она всё ещё практична. Некоторые Android-телефоны (например, отдельные модели Xiaomi и Huawei) оснащены встроенным ИК-передатчиком и могут управлять телевизорами, кондиционерами, ТВ-приставками и другой бытовой техникой.

\booksubsection{Принцип работы:}

ИК-пульт излучает через инфракрасный светодиод импульсные сигналы определённой частоты (обычно 38 к Гц), а ИК-приёмник бытовой техники декодирует эти импульсы и выполняет соответствующую команду. Разные бренды используют разные протоколы кодирования, наиболее распространённые: NEC, RC5, RC6, SONY и др. Приложение ИК-пульта в телефоне содержит встроенную библиотеку кодов разных брендов — выберите бренд и модель, и оно подберёт нужный код.

\booksubsection{Ограничения:}

\begin{itemize}
  \item \textbf{Направленность}: инфракрасный луч распространяется по прямой, нужно направлять на устройство, не проходит сквозь стены.
  \item \textbf{Малое расстояние}: эффективная дальность обычно 5–8 метров.
  \item \textbf{Apple не поддерживает}: у iPhone нет ИК-передатчика, эта функция недоступна.
\end{itemize}

\booksection{NFC: магия одного касания}

NFC (Near Field Communication, ближняя бесконтактная связь) — это технология высокочастотной беспроводной связи на сверхкоротком расстоянии (до 10 см), рабочая частота 13,56 МГц, максимальная скорость передачи 424 Кбит/с.

\booksubsection{Принцип связи и режимы:}

NFC развился из технологии RFID и поддерживает три режима работы:

\begin{itemize}
  \item \textbf{Режим чтения/записи}: телефон считывает информацию с NFC-меток, например, считывает NFC-стикер для запуска действия.
  \item \textbf{Режим точка-точка}: два NFC-устройства обмениваются данными, например Android Beam (уже отключён).
  \item \textbf{Режим эмуляции карты}: телефон эмулирует смарт-карту — это самый используемый режим: транспортные карты, карты доступа, банковские карты работают именно в нём.
\end{itemize}

\booksubsection{Механизмы безопасности:}

\begin{itemize}
  \item \textbf{Сверхкороткое расстояние}: дальность связи до 10 см сама по себе является физической защитой, удалённый перехват практически невозможен.
  \item \textbf{Защищённый чип}: телефон содержит независимый защищённый элемент (SE) или использует HCE (Host Card Emulation), платёжные ключи хранятся в защищённой области, их невозможно извлечь даже при root-доступе.
  \item \textbf{Подтверждение по отпечатку/лицу}: при оплате требуется биометрическое подтверждение, даже при утере телефона мошенническая оплата невозможна.
  \item \textbf{Токенизация}: Apple Pay, Huawei Pay и другие не передают реальный номер карты, а используют одноразовый токен — даже перехваченный, он не может быть использован повторно.
\end{itemize}

\booksubsection{Типичные применения:}

\begin{itemize}
  \item \textbf{Транспортная карта / карта доступа}: запишите данные физической карты в телефон и прикладывайте телефон.
  \item \textbf{Мобильные платежи}: Apple Pay, Huawei Pay и др. — привяжите банковскую карту к телефону, платите одним касанием.
  \item \textbf{Сопряжение устройств}: быстрое сопряжение Bluetooth-наушников и колонок касанием NFC.
  \item \textbf{Проверка личности}: считывание NFC-удостоверений личности, проверка электронных водительских прав.
\end{itemize}

\booksection{Персональная точка доступа: превращаем телефон в Wi-Fi-роутер}

Когда у вас много трафика, можно раздать его на планшет или ноутбук.

\begin{itemize}
  \item \textbf{Как включить}: Настройки > Персональная точка доступа / Общий доступ к сети, включите её.
  \item \textbf{На заметку}: сильно расходует батарею и трафик, обязательно установите сложный пароль и ограничьте количество подключённых устройств.
  \item \textbf{USB-модем}: подключите телефон к компьютеру кабелем — стабильнее и быстрее, чем беспроводная точка доступа.
\end{itemize}

\begin{tcbpractice}{Попробуйте сами: настройка и управление точкой доступа}
\begin{enumerate}
  \item Откройте [Настройки], найдите «Персональная точка доступа» или «Общий доступ к сети», включите, задайте имя и пароль.
  \item С другого устройства (например, ноутбука) найдите и подключитесь к этому Wi-Fi.
  \item Посмотрите на телефоне, можно ли управлять подключёнными устройствами или отключить незнакомые.
\end{enumerate}

\end{tcbpractice}
